\subsection*{Terminology}
\addcontentsline{toc}{subsection}{Terminology}
\label{sec:front-definitions}

Table~\ref{tab:definitions} collects the symbols and conventions used throughout, each with the place that defines it, and the names the note uses for the four components of the analysis chain, with the code identifiers it cites.

\begin{table}[H]
\centering\small
\caption{Definitions and conventions, each with the section or table that defines it, and the components of the analysis chain with their code identifiers.}
\label{tab:definitions}
\begin{tabular}{@{}>{\raggedright\arraybackslash}p{0.2\textwidth}>{\raggedright\arraybackslash}p{0.76\textwidth}@{}}
\toprule
Term & Definition \\
\midrule
$\pT$, $\pPar$ & Muon momentum transverse and parallel to the beam axis, in GeV/$c$ (Table~\ref{tab:observables-impl}). \\
$\Eavail$ & Available energy: proton kinetic energy, charged-pion energy less \SI{135}{MeV}, and $\pi^{0}$ and $\gamma$ energy, in GeV (\S\ref{sec:eavail-definition}). \\
$q_{3}$ & Three-momentum transfer, in GeV/$c$ (Table~\ref{tab:observables-impl}). \\
$W$ & Hadronic invariant mass, in GeV. Truth $W$ uses the struck-nucleon mass, reconstructed $W$ the average nucleon mass. The external generator predictions use the neutron mass for every event (\S\ref{sec:eavailw-axis}). \\
Data & The twelve ME FHC playlists, $1.057\times10^{21}$ protons on target (\S\ref{sec:experiment}). \\
Abbreviations & IBU: D'Agostini iterative Bayesian unfolding (\S\ref{sec:intro}). ME, FHC: the medium-energy, neutrino-mode (forward horn current) exposure (\S\ref{sec:experiment}). POT: protons on target (App.~\ref{app:codebase}). GBDT: gradient-boosted decision tree (\S\ref{sec:intro}). PPFX: the hadron-production weights of the flux prediction (\S\ref{sec:experiment}). MAT: the MINERvA Analysis Toolkit (\S\ref{sec:experiment}). \\
\addlinespace
Configurations & (a)--(f): Table~\ref{tab:config-key}, from Table~\ref{tab:estimators-matrix} and \S\ref{sec:method-config}. \\
\addlinespace
$Q_R$ & Covariance of the two-dimensional cross section: the sum of the training-seed, statistical, systematic-universe and run-period blocks (\S\ref{sec:systematics}, Table~\ref{tab:uq}). \\
$C_{\mathrm{total}}+C_{\mathrm{resp}}+C_{\mathrm{PL,5D}}$ & The reported covariance of the sixteen cells, with $C_{\mathrm{total}}=C_{\mathrm{e}}+C_{\mathrm{model}}$ (\S\ref{sec:fived16-cov}). \\
$C_{\mathrm{e}}$ & Sums the statistical bootstrap of configuration~(e), the PPFX flux universes, the interaction-model and detector bands, the lateral bands, the normalization and the training-seed block (\S\ref{sec:fived16-cov}). \\
$C_{\mathrm{model}}$ & The unscaled outer-product sum of the data result's shifts when the prior is reweighted in turn by the ratio of GENIE + MEC, NuWro or GiBUU to GENIE CV, in cross-section units. It measures generator-shape model dependence and is not a calibrated bound (\S\ref{sec:fived16-cov}). \\
$C_{\mathrm{resp}}$ & The sum, over four deposit categories (protons, mesons, electromagnetic and other deposits), of the outer product of the half-difference of the two unfolded shifts when that category's reconstructed energy is scaled by its one-standard-deviation factor (\S\ref{sec:fived16-cov}). \\
$C_{\mathrm{PL}}$, $C_{\mathrm{PL,5D}}$ & The run-period blocks of the two-dimensional result and of the sixteen cells, built from the double ratios of the run-period groups' unfolded results with no reference period (\S\ref{sec:uq-combine}, \S\ref{sec:fived16-pl}). \\
$V_g$ & A generator prediction's finite-sample (diagonal) covariance (Table~\ref{tab:fived16-gen}). \\
$C_Z$ & Five-axis covariance of the configuration-(a) central value. It is not an uncertainty product and not the sixteen-cell covariance; \S\ref{sec:5d} uses it only as a descriptive scale (\S\ref{sec:eavailw-adopted}). \\
$C_k$ & Coverage within $\pm k$ standard deviations (\S\ref{sec:limitations}). \\
\addlinespace
Cells & R00--R15: the sixteen reported cells, formed from the H2 grid by purity merges; each cell's box is in Table~\ref{tab:fived16-cells} (\S\ref{sec:fived16}). The 42 $(\Eavail,W)$ cells and the \jtNcells{} coarse joint cells of the $3^{5}$ grid (Fig.~\ref{fig:bkg-residual-note}; \S\ref{sec:joint5d}). \\
J, H2, EW & Reporting grids of the five-dimensional studies. J: 109 supported joint cells subdividing all five axes. H2: 27 supported coarser joint cells, also subdividing all five axes, not nested within J. EW: 39 supported $(\Eavail,W)$ cells (42 with the three unsupported). Boundaries and the support rule: App.~\ref{sec:gbdt-diagnostic-definitions}; labels: \S\ref{sec:gbdt-model-dependence}. \\
D1, W1--W3 & Fixed historical truth deformations, built before the generator flux was repaired: D1 a GiBUU/GENIE $\Eavail$ ratio, W1 a NuWro/GENIE $\Eavail$ ratio, W2 a GENIE MEC/CV $(\Eavail,W)$ ratio, W3 a NuWro/GENIE $(\pT,\pPar,\Eavail)$ ratio (\S\ref{sec:gbdt-model-dependence}). \\
$R$ & The raised-refinement-capacity estimator of \S\ref{sec:val-departure}: baseline OmniFold classifiers with a 400-tree, 31-leaf background-refinement classifier (\S\ref{sec:gbdt-model-dependence}). \\
RD1--RD3 & Reporting definitions of the five-dimensional precision study at J, H2 and EW resolution (\S\ref{sec:gbdt-model-dependence}). \\
Units & Densities in \si{cm^2/nucleon} divided by the product of the kept axes' units; cell-integrated cross sections in \si{cm^2/nucleon} (App.~\ref{app:release}, Table~\ref{tab:fived16-cells}). \\
\addlinespace
Event loop (C++) & \texttt{runEventLoopOmniFold}, built on the MINERvA Analysis Toolkit (MAT) and the MINERvA-101 framework. It reads production AnaTuples, applies selection and systematic-universe machinery, and writes per-event ROOT trees. Code: \repopath{MINERvA101/MINERvA-101-Cross-Section/runEventLoopOmniFold.cpp}; its four trees are \texttt{mc\_truth\_denom} (efficiency denominator), \texttt{mc\_signal\_reco}, \texttt{mc\_background} and \texttt{data}; the flux histogram is \texttt{pTmu\_reweightedflux\_integrated}. \\
Event-loop outputs & Called \emph{omnifiles} in the code: per-event records for data, for simulated signal with paired true and reconstructed quantities, for simulated background, and for the truth-level efficiency denominator, together with the exposure in protons on target (POT). The \texttt{\_universes\_full} variants add per-universe weight and shifted-kinematics branches. \\
Driver (Python) & Per-analysis orchestration, by a two-dimensional driver and its $N$-dimensional generalization. The driver builds training arrays from the event-loop outputs, subtracts the background by the method an option selects, the purity weight by default (\S\ref{sec:method-algo}), or negative-weight injection with or without the Stay-Positive refinement (App.~\ref{app:negweight}), runs the engine, and extracts the cross section. Cross-section extraction (\S\ref{sec:method-xsec}): \texttt{extract\_cross\_section\_2d} in the 2D driver, \texttt{2d-unfolding/unfold\_2d\_omnifold\_unbinned.py}, and its N-dimensional generalization \texttt{nd-unfolding/xsec\_nd.py}. \\
        % Accuracy fix (unscored; director ruling 2026-10-26): the option is --bkg-mode {purity, negweight,
        % negweight-refined}, default purity (2d-unfolding/unfold_2d_omnifold_unbinned.py:1362-1377 and :1825-1870 @ ddf16f10; nd-unfolding/unfold_nd_omnifold_unbinned.py:619-630 and :900-1030 @ ddf16f10);
        % the 2D result uses the refined negative-weight mode (sec:negweight-footing).
Engine (Python) & \texttt{unbinned\_unfolding/python/omnifold.py}: the OmniFold iteration, implemented independently of the learner backend: LightGBM for the central value and the uncertainty ensembles, scikit-learn's exact-split gradient boosting (\texttt{exact}) for the cross-check, its histogram gradient boosting (\texttt{hist}) as an alternate, and a multilayer perceptron (MLP) for a further cross-check. \\
Figure code & For the figures of \S\ref{sec:results}, \texttt{2d-unfolding/plot\_2d\_paper\_comparison.py} builds the simulated-truth curve of Fig.~\ref{fig:ptslices} from \texttt{hOFTruthDenom2D}; \texttt{2d-unfolding/compare\_to\_models.py} draws Fig.~\ref{fig:modelmaps}. \\
\bottomrule
\end{tabular}
\end{table}


